\subsection{Постановка задачи}

Микропроект посвящён двумерному моделированию фазового перехода лёд--вода
(задача Стефана) энтальпийным методом конечных разностей.
Задача Стефана~--- это задача теплопроводности с подвижной границей:
при плавлении или замерзании граница между фазами перемещается,
поглощая или выделяя скрытую теплоту $L$.

Классическая формулировка требует отслеживания фронта фазового перехода
и выполнения на нём условия Стефана:
\begin{equation}\label{eq:stefan_condition}
  \rho L \frac{ds}{dt} = k_l \frac{\partial T_l}{\partial n}\bigg|_{s}
                       - k_s \frac{\partial T_s}{\partial n}\bigg|_{s},
\end{equation}
где $s(t)$~--- положение фронта, $k_l$, $k_s$~--- теплопроводности жидкой
и твёрдой фаз, $\partial T / \partial n$~--- производные температуры по
нормали к границе с каждой стороны. Это условие выражает баланс тепловых
потоков: разница потоков идёт на плавление (или намерзание) вещества.

В двумерном случае явно отслеживать фронт уже сложно.
Энтальпийный метод позволяет обойти эту проблему, заменяя явное условие
на границе неявной обработкой через функцию $T(H)$.

\subsection{Математическая модель}

\subsubsection{Физические основы}

Рассмотрим среду, состоящую из одного вещества (вода/лёд) с фазовым переходом
I рода при температуре $T_{\rm melt}$. При фазовом переходе I рода
выделяется или поглощается скрытая теплота $L$ (Дж/кг) --- энергия,
необходимая для разрыва межмолекулярных связей кристаллической решётки льда,
не изменяющая при этом температуру вещества.

Теплоперенос в каждой фазе описывается уравнением Фурье:
\begin{equation}\label{eq:fourier}
  \rho c \frac{\partial T}{\partial t} = k\,\nabla^2 T,
\end{equation}
где $\rho$~--- плотность, $c$~--- удельная теплоёмкость, $k$~--- коэффициент
теплопроводности. Это уравнение следует из закона сохранения энергии
и закона Фурье $\mathbf{q} = -k\,\nabla T$, связывающего тепловой поток
с градиентом температуры.

Параметры вещества различаются для твёрдой и жидкой фаз:

\begin{center}
\begin{tabular}{lcc}
  \hline
  Параметр & Лёд & Вода \\
  \hline
  Теплопроводность $k$, Вт/(м$\cdot$К) & 2.2 & 0.6 \\
  Удельная теплоёмкость $c$, Дж/(кг$\cdot$К) & 2090 & 4180 \\
  Скрытая теплота $L$, Дж/кг & \multicolumn{2}{c}{334\,000} \\
  Плотность $\rho$, кг/м$^3$ & \multicolumn{2}{c}{1000} \\
  \hline
\end{tabular}
\end{center}

Температуропроводность $\alpha = k / (\rho c)$ определяет скорость
диффузии тепла: для льда $\alpha_s \approx 1{,}05 \cdot 10^{-6}$~м$^2$/с,
для воды $\alpha_l \approx 1{,}44 \cdot 10^{-7}$~м$^2$/с. Лёд проводит
тепло почти на порядок быстрее воды.

\subsubsection{Энтальпийная формулировка}

Введём объёмную удельную энтальпию $H$ (Дж/м$^3$), отсчитанную от точки
плавления. В этой задаче энтальпия показывает, сколько тепловой энергии
накоплено в единице объёма вещества. Уравнение
теплопроводности в терминах энтальпии принимает единую форму для обеих фаз:
\begin{equation}\label{eq:enthalpy_pde}
  \frac{\partial H}{\partial t} = k\,\nabla^2 T.
\end{equation}

Связь температуры и энтальпии кусочно-линейна:
\begin{equation}\label{eq:T_of_H}
  T(H) = \begin{cases}
    T_{\rm melt} + \dfrac{H}{\rho c_s},          & H < 0          \quad\text{(лёд)}, \\[8pt]
    T_{\rm melt},                                  & 0 \le H \le \rho L \quad\text{(фазовый переход)}, \\[8pt]
    T_{\rm melt} + \dfrac{H - \rho L}{\rho c_l},  & H > \rho L     \quad\text{(вода)}.
  \end{cases}
\end{equation}

Физический смысл трёх областей:
\begin{itemize}
  \item $H < 0$: вещество полностью в твёрдой фазе; вся энергия расходуется
        на изменение температуры;
  \item $0 \le H \le \rho L$: вещество находится в процессе фазового перехода;
        подводимая энергия уходит на фазовый переход, температура при
        этом остаётся равной $T_{\rm melt}$;
  \item $H > \rho L$: вещество полностью в жидкой фазе; вся скрытая теплота
        уже поглощена, дальнейшая энергия снова идёт на нагрев.
\end{itemize}

Ключевое достоинство энтальпийного метода: граница раздела фаз
обрабатывается автоматически. Ячейки, в которых $0 \le H \le \rho L$,
находятся <<в процессе>> фазового перехода~--- фронт плавления
распределяется по нескольким ячейкам сетки, не требуя явного
отслеживания его положения.

\subsubsection{Теплопроводность в зоне фазового перехода}

В ячейках, находящихся в процессе плавления ($0 \le H \le \rho L$),
теплопроводность интерполируется между фазами:
\begin{equation}
  k(H) = (1 - f)\,k_s + f\,k_l, \qquad f = \frac{H}{\rho L},
\end{equation}
где $f \in [0,1]$~--- доля жидкой фазы. Это обеспечивает плавный
переход теплофизических свойств через фронт.

\subsection{Численный метод}

\subsubsection{Явная конечно-разностная схема}

Уравнение~(10) дискретизируется на равномерной
квадратной сетке $N \times N$ с шагом $h = D / N$, где $D$~--- размер
области. Лапласиан аппроксимируется стандартным пятиточечным шаблоном:
\begin{equation}\label{eq:scheme}
  H_{i,j}^{n+1} = H_{i,j}^{n}
    + \Delta t \cdot k_{i,j}
    \cdot \frac{T_{i+1,j}^n + T_{i-1,j}^n + T_{i,j+1}^n + T_{i,j-1}^n
          - 4\,T_{i,j}^n}{h^2},
\end{equation}
где $k_{i,j} = k(H_{i,j}^n)$~--- теплопроводность, зависящая от фазы.

Схема является \textbf{явной}: значения на новом временном слое $n+1$
вычисляются непосредственно из значений на слое $n$, без решения
системы линейных уравнений. Это даёт простую реализацию, но накладывает
ограничение на шаг по времени.

\subsubsection{Условие устойчивости}

Явная схема устойчива при выполнении условия CFL:
\begin{equation}\label{eq:cfl}
  \Delta t \le \frac{h^2\,\rho\,c_{\min}}{4\,k_{\max}},
\end{equation}
где множитель 4 соответствует двумерному случаю (четыре соседа в шаблоне),
$c_{\min}$ и $k_{\max}$~--- наименьшая теплоёмкость и наибольшая
теплопроводность среди фаз. Физический смысл: за один шаг по времени
тепло не должно <<перескочить>> через ячейку.

Для используемых параметров ($h = 10^{-4}$~м, $k_{\max} = 2{,}2$,
$c_{\min} = 2090$, $\rho = 1000$):
\[
  \Delta t_{\max} \approx 2{,}4 \cdot 10^{-3}~\text{с}.
\]

\subsubsection{Граничные условия}

Реализованы два типа граничных условий:

\textbf{Фиксированные (Дирихле):}
температура на границе не меняется~--- $T|_{\partial\Omega} = T_0$.
Используется для моделирования контакта с телом заданной температуры
(например, рука при $36\,^{\circ}$C).

\textbf{Адиабатические (Нейман):}
тепловой поток через границу равен нулю~---
$\partial T / \partial n|_{\partial\Omega} = 0$.
Реализуется копированием энтальпии из ближайшей внутренней ячейки.
Используется для замкнутых систем без теплообмена с окружением.

\subsubsection{Двойная буферизация}

При обновлении сетки по схеме~(13) значения $T_{i\pm1,j}^n$
и $T_{i,j\pm1}^n$ должны быть со \textit{старого} временного слоя.
Поэтому используется буфер $H_{\rm new}$: сначала все новые значения
записываются в буфер, затем буфер копируется обратно в основной массив.

\subsection{Валидация}

\subsubsection{Одномерная задача Стефана}

Для верификации численной схемы рассмотрена одномерная задача: полупространство
льда при $T = T_{\rm melt}$, левая граница мгновенно нагревается до $T_{\rm hot}$.
Аналитическое решение даёт положение фронта плавления:
\begin{equation}\label{eq:stefan_exact}
  s(t) = 2\lambda\sqrt{\alpha_l\, t},
\end{equation}
где $\alpha_l = k_l / (\rho c_l)$~--- температуропроводность жидкой фазы,
а $\lambda$~--- корень трансцендентного уравнения:
\begin{equation}\label{eq:transcendental}
  \mathrm{St} = \lambda\sqrt{\pi}\,e^{\lambda^2}\,\mathrm{erf}(\lambda),
\end{equation}
где $\mathrm{St} = c_l (T_{\rm hot} - T_{\rm melt}) / L$~--- число Стефана,
характеризующее отношение теплоты, затраченной на нагрев жидкости,
к скрытой теплоте плавления. При $T_{\rm hot} = 50\,^{\circ}$C:
\[
  \mathrm{St} = \frac{4180 \cdot 50}{334\,000} \approx 0{,}626.
\]

\begin{figure}[H]
  \centering
  \includegraphics[width=1.00\linewidth]{img/micro/validation.png}
  \caption{Сравнение численного положения фронта $s(t)$ с аналитическим
           решением. Численный фронт
           воспроизводит закон $s \propto \sqrt{t}$.}
\end{figure}

\subsection{Постановки задач (сцены)}

\subsubsection{Таяние снежинки}

Снежинка задаётся PNG-маской: тёмные пиксели маски инициализируются как лёд
($H = \rho c_s \cdot (T_0 - T_{\rm melt}) < 0$), окружающая среда~--- как
тёплая вода ($H = \rho L + \rho c_l \cdot (T_w - T_{\rm melt}) > \rho L$).
Тонкие лучи тают первыми, так как у них наибольшее отношение поверхности
к объёму: тепло проникает к центру тонкого элемента быстрее.

\subsubsection{Таяние снежинки на руке}

Снежинка помещается на нижнюю границу области с фиксированной температурой
$T_{\rm hand} = 36\,^{\circ}$C (температура тела человека). Окружающий воздух
моделируется как лёд при $T \approx 0\,^{\circ}$C. При размере снежинки $\sim 5$~мм
и реальных теплофизических параметрах время полного таяния составляет 2--3~с,
что качественно совпадает с повседневным опытом.

\subsubsection{Намерзание (сосулька)}

Вода при $+2\,^{\circ}$C, верхняя стенка охлаждена до $-20\,^{\circ}$C,
боковые стенки при $+5\,^{\circ}$C. Лёд нарастает сверху вниз; тёплые
боковые границы не дают ему расшириться, формируя вытянутую структуру.
Эта постановка демонстрирует обратимость фазового перехода: тот же
солвер без изменений моделирует и плавление, и замерзание.

\subsection{Реализация}

\subsubsection{Архитектура}

Проект реализован на C++17 с разделением на модули:
\begin{itemize}
  \item \texttt{core/}~--- физика (Grid, Solver, Material). Не зависит от
        библиотеки визуализации;
  \item \texttt{render/}~--- визуализация (Renderer, Colormap). Зависит
        от raylib, не знает о деталях физики;
  \item \texttt{scenes/}~--- начальные условия. Каждая сцена реализует
        интерфейс \texttt{Scene};
  \item \texttt{export/}~--- вывод данных в CSV для анализа в Python.
\end{itemize}

Зависимости однонаправлены: \texttt{render} $\to$ \texttt{core} (только чтение),
\texttt{scenes} $\to$ \texttt{core}. Циклических зависимостей нет.

\subsubsection{Оптимизация рендеринга}

Наивная отрисовка $N \times N$ прямоугольников через \texttt{DrawRectangle}
($\sim 40\,000$ вызовов при $N = 200$) даёт $\sim 4$~FPS. Вместо этого
температурное поле записывается в массив пикселей, загружается как текстура
на GPU (\texttt{UpdateTexture}) и отрисовывается одним вызовом
\texttt{DrawTextureEx}. Результат: $\sim 60$~FPS.

\subsubsection{Интерактивное управление}

Raylib-приложение поддерживает:
\begin{itemize}
  \item \texttt{[1]--[5]}~--- переключение сцен;
  \item \texttt{[SPACE]}~--- пауза / продолжение;
  \item \texttt{[R]}~--- сброс состояния;
  \item \texttt{[LMB]}~--- локальный нагрев мышью;
  \item \texttt{[RMB]}~--- локальное охлаждение мышью;
  \item \texttt{[F]}~--- запись кадров в PNG.
\end{itemize}

\subsection{Результаты}



\begin{figure}[H]
  \centering
  \includegraphics[width=0.8\linewidth]{img/micro/image.png}
    \caption{Таяние снежинки. Синий~--- лёд, белый~--- фазовый переход,
           красный~--- вода. Тонкие лучи исчезают первыми.}
\end{figure}

См. \href{https://www.youtube.com/watch?v=3Jm9fgB6ytE&feature=youtu.be}{Видео на YouTube}.

Качественные наблюдения:
\begin{enumerate}
  \item Тонкие элементы снежинки тают быстрее массивных~--- следствие
        большего отношения площади поверхности к объёму.
  \item Фронт плавления в 1D воспроизводит аналитический закон
        $s \propto \sqrt{t}$.
  \item Намерзание и плавление обрабатываются одним и тем же кодом без
        изменений~--- обратимость фазового перехода I рода.
\end{enumerate}

\subsection{Выводы}

Реализован двумерный решатель задачи Стефана энтальпийным методом
на явной конечно-разностной схеме. Энтальпийная формулировка позволяет
избежать явного отслеживания фронта плавления, сохраняя физически
корректную картину фазового перехода. Валидация на одномерном
аналитическом решении подтвердила корректность реализации.

Продемонстрированы три режима: таяние (снежинка в тёплой воде),
контактное плавление (снежинка на руке) и намерзание (рост льда
от холодного источника). Все режимы обрабатываются единым солвером
без изменения уравнений~--- это следствие энтальпийной
формулировки.
